% Lösungen Einheit 4 von 4 – Sachaufgaben (zusammenbau v0.1)
\einheitenkopf{Einheit 4 von 4 · Sachaufgaben}

%% TODO Lösung der Erklärzeile 21g) fehlt in der Bank
\erg{21}{a) beide – die Frage schließt negative Zahlen nicht aus \quad b) $x_1 = 12$, $x_2 = -12$ – beide Zahlen passen \quad c) $x_1 = 20$, $x_2 = -20$ – beide Zahlen passen \quad d) $x_1 = 50$, $x_2 = -50$ – beide Zahlen passen \quad e) $x_1 = 18$, $x_2 = -18$ – beide Zahlen passen \quad f) $x_1 = 25$, $x_2 = -25$ – beide Zahlen passen \quad g) \ldots \quad h) $x^2 + 7 = 88$; $x^2 = 81$; $x_1 = 9$, $x_2 = -9$ \quad i) $x \cdot (x + 1) = 56$; $x^2 + x - 56 = 0$; $x_{1,2} = -0{,}5 \pm \sqrt{0{,}25 + 56} = -0{,}5 \pm 7{,}5$; $x_1 = 7$, $x_2 = -8$; Zahlenpaare $7$ und $8$ oder $-8$ und $-7$ \quad j) $x \cdot (x + 5) = 84$; $x^2 + 5x - 84 = 0$; $x_{1,2} = -2{,}5 \pm \sqrt{6{,}25 + 84} = -2{,}5 \pm 9{,}5$; $x_1 = 7$, $x_2 = -12$ \quad k) $x \cdot (x + 4) = 96$; $x^2 + 4x - 96 = 0$; $x_{1,2} = -2 \pm \sqrt{4 + 96} = -2 \pm 10$; $x_1 = 8$, $x_2 = -12$; Breite $8$ cm – $-12$ fällt weg, eine Länge ist nie negativ \quad l) $x \cdot (x + 40) = 6000$; $x^2 + 40x - 6000 = 0$; $x_{1,2} = -20 \pm \sqrt{400 + 6000} = -20 \pm 80$; $x_1 = 60$, $x_2 = -100$; Länge $= 60 + 40 = 100$ cm. Antwort: Es ist $60$ cm breit und $100$ cm lang. \quad m) Länge $7 + 3 = 10$ cm; $7 \cdot 10 = 70$ (wA): stimmt \quad n) $(x + 2) \cdot (x + 6) = 96$; $x^2 + 8x + 12 = 96$; $x^2 + 8x - 84 = 0$; $x_{1,2} = -4 \pm \sqrt{16 + 84} = -4 \pm 10$; $x_1 = 6$, $x_2 = -14$; $-14$ fällt weg: Breite $= 6$ m, Länge $= 10$ m}
\erg{22}{a) $(x + 10)^2 = 900$; $x + 10 = 30$ oder $x + 10 = -30$; $x_1 = 20$, $x_2 = -40$; Seite $20$ cm}
\erg{23}{a) Eine Länge ist nie negativ, $-7$ fällt weg: Breite $= 5$ m, Länge $= 7$ m \quad b) Eine Breite ist eine Länge und nie negativ; $-13$ m gibt es nicht}
